\documentclass[12pt]{amsart}
\usepackage{amsmath,amsthm,amssymb}
\usepackage{dsfont,bm,mathtools}
\usepackage[usenames,dvipsnames,svgnames,table]{xcolor}
\usepackage[top=2.3cm,bottom=2.3cm,left=2cm,right=2cm]{geometry}
\usepackage{hyperref}
\hypersetup{hidelinks}
\allowdisplaybreaks
\emergencystretch=1.5em
\definecolor{ProposalOrange}{RGB}{160,65,0}
\definecolor{RefereePurple}{RGB}{103,34,147}
\newtheorem{thm}{Theorem}[section]
\newtheorem{lem}[thm]{Lemma}
\theoremstyle{remark}
\newtheorem{rem}[thm]{Remark}
\theoremstyle{definition}
\newtheorem{nota}[thm]{Notation}
\definecolor{ResolutionGreen}{RGB}{0,100,55}
\newcommand{\C}{\mathds{C}}
\newcommand{\Q}{\mathds{Q}}
\newcommand{\rv}{\operatorname{rv}}
\newcommand{\Gal}{\operatorname{Gal}}
\newcommand{\Aut}{\operatorname{Aut}}
\newcommand{\RV}{\mathrm{RV}}
\newcommand{\RES}{\mathrm{RES}}
\newcommand{\VF}{\mathrm{VF}}
\newcommand{\K}{\mathds{k}}
\newcommand{\Z}{\mathds{Z}}
\newcommand{\N}{\mathds{N}}
\newcommand{\vrv}{\operatorname{vrv}}
\newcommand{\GL}{\operatorname{GL}}
\newcommand{\dcl}{\operatorname{dcl}}
\newcommand{\pr}{\operatorname{pr}}
\newcommand{\defn}{\mathrm{def}}
\newcommand{\given}{\,\middle|\,}
\newcommand{\Sbase}{\mathbb S}
\newcommand{\RVlang}{\mathcal L_{\mathrm{RV}}}
\numberwithin{equation}{section}
\pagestyle{plain}

\newcommand{\A}{\mathds{A}}
\newcommand{\Gm}{\mathbb G_m}
\newcommand{\Var}{\mathrm{Var}}
\newcommand{\Kres}{K_0(\RES)}
\newcommand{\Bangres}{{!K_0(\RES)}}
\newcommand{\Kraw}{K_0^{\widehat\mu,\mathrm{raw}}(\Var_{\C})}
\newcommand{\Kmon}{K_0^{\widehat\mu}(\Var_{\C})}
\newcommand{\E}{\mathbb E_b}
\newcommand{\Vol}{\operatorname{Vol}}
\newcommand{\vv}{\operatorname{val}}
\newcommand{\LiftingTheorem}{the lifting theorem of Section~3}

\theoremstyle{plain}
\newtheorem*{thm*}{Theorem}

\newcommand{\R}{\mathds R}
\newcommand{\Sfiber}{\mathbb S_b}
\newcommand{\Rbase}[1]{{!K_0(\RES_{#1})}}
\newcommand{\IsVFVG}{(\VF,\Gamma)\text{-generated}}
\newcommand{\gsk}{K_+}
\theoremstyle{plain}
\newtheorem{test}[thm]{Test}

\theoremstyle{plain}
\newtheorem{cor}[thm]{Corollary}
\begin{document}
\setcounter{section}{5}
\setcounter{thm}{20}
{\small\noindent Forey--Yin, \emph{Bounded integral and motivic Milnor fiber}.
Section 5: remaining corollary.
Approved context is black; proposed manuscript text is orange.}
\medskip

\noindent\emph{End of the proof of Theorem 5.17.}
The finite witnesses give equality of the global bounded
retractions in $R(\Sbase)$.
All cancellations take place in its additive group;
no torus class is inverted.
For $F=\C((t))$, the comparison $\Theta$ transports this
global equality to motivic volume.

{\color{ProposalOrange}
\begin{cor}\label{cor:global:vanishing}
Let $F=\C((t))$ and let $f:A\to B$ be a morphism of
varieties over $\C$.
Interpret these varieties in a sufficiently saturated
ambient model $\mathbb U$ of $\mathrm{ACVF}(0,0)$, using
finite affine charts.
For $b\in B(\mathbb U)$ put $A_b=f^{-1}(b)$ and
$\Sfiber=F\langle b\rangle$, with $b$ represented by its
VF chart coordinates.
If, for every $b\in B(\mathbb U)$,
\[
 \mathbb E_{\mathrm b}^{\Sfiber}\!
       \left(\int[A_b]\right)=0
       \quad\text{in }R(\Sfiber),
\]
then $\Vol([A])=0$ in $\Kmon$.
\end{cor}

\begin{proof}
Choose finite disjoint constructible pieces subordinate to
affine covers of $A$ and $B$ and encode their finite tags
by distinct elements of $\C$.
This realizes $f$ as a definable map between VF tuples.
Changing charts gives definably equivalent parameter
bases and transports the fiber equality.

Both $F$ and each $\Sfiber$ in these coordinates are
VF-generated. Thus the integrals of Section~3 are
available over these bases.
The lifting theorem supplies positive RV representatives
of $A_b$ with actual $\Sfiber$-definable lift bijections.
Apply Theorem~5.17 to $A\to B$ and the empty family
$\varnothing\to B$.
The displayed hypothesis is precisely the required
equality of their fiber retractions, so
\[
 \mathbb E_{\mathrm b}^{F}\!
       \left(\int[A]\right)=0
       \quad\text{in }R(F).
\]
Applying the global comparison $\Theta$ gives
$\Vol([A])=0$.
\end{proof}
}

\medskip
{\small\noindent\color{ProposalOrange}
\textbf{Proposed disposition.}
Replace the last corollary by this statement and proof.
The fiber hypothesis uses the intrinsic coefficient ring
already retained in Remark~5.18; it requires no comparison
$\Theta_b$ over arbitrary fiber bases.
The quantifier includes all ambient points, as required
by Theorem~5.17, rather than only Puiseux points.
This is the zero-fiber specialization of the corrected
Fubini theorem.}

\end{document}
